\subsection{Variational characterizations of eigenvalues}

In this section, we assume that K = C.

Let \(u\) be a compact and positive endomorphism of E and let \((\lambda_n(u))_{n\in \mathrm{I_E}}\) be the decreasing sequence of eigenvalues of \(u\) . We set \(\mathrm{I} = \mathrm{I_E}\) .

For every closed subspace F of E, we denote by

\[
r _ {\mathrm{F}} (u) = \inf _ {x \in \mathrm{F} - \{0 \}} \frac {\langle x \mid u (x) \rangle}{\| x \| ^ {2}}, \quad \mathrm{R} _ {\mathrm{F}} (u) = \sup _ {x \in \mathrm{F} ^ {\circ} - \{0 \}} \frac {\langle x \mid u (x) \rangle}{\| x \| ^ {2}},
\]

where the lower bound (resp. the upper bound) is taken in \([0, +\infty]\) .

For every \(n \in \mathbf{N}\) , we denote by \(\mathcal{F}_n\) the set of vector subspaces \(\mathrm{F} \subset \mathrm{E}\) of dimension \(n\) . We say that a subspace \(\mathrm{F} \in \mathcal{F}_n\) is adapted to \(u\) if it admits an orthonormal basis \((f_i)_{0 \leqslant i \leqslant n-1}\) such that \(u(f_i) = \lambda_i(u)f_i\) for \(0 \leqslant i \leqslant n-1\) .

PROPOSITION 5. --- Let \(n \in I\) .

\begin{enumerate}
\def\labelenumi{\alph{enumi})}
\item
  For every subspace \(\mathrm{F} \in \mathcal{F}_{n+1}\) adapted to \(u\) , we have \(\lambda_n(u) = r_{\mathrm{F}}(u)\) ;
\item
  For every subspace \(F \in F_{n}\) adapted to u, we have \(\lambda_{n}(u) = R_{\mathrm{F}}(u)\) .
\end{enumerate}

Let \(F \in F_{n+1}\) be a subspace adapted to u, and \((f_i)_{0 \leqslant i \leqslant n}\) be an orthonormal basis of F such that \(u(f_i) = \lambda_i(u)f_i\) for all i. For all x in F, we have

\[
\langle x \mid u (x) \rangle = \sum_ {0 \leqslant i \leqslant n} \lambda_ {i} (u) | \langle f _ {i} \mid x \rangle | ^ {2} \geqslant \lambda_ {n} (u) \sum_ {0 \leqslant i \leqslant n} | \langle f _ {i} \mid x \rangle | ^ {2} = \lambda_ {n} (u) \| x \| ^ {2},
\]

with equality if \(x = f_{n}\) . This implies that \(r_{\mathrm{F}}(u) = \lambda_{n}(u)\) whence assertion (a).

Let \(F \in F_{n}\) be a subspace adapted to u. The closed subspace \(F^{\circ}\) is nonzero (since \(n = \dim(F) < \dim(E)\) ) and stable under u; the endomorphism \(\widetilde{u}\) of \(F^{\circ}\) deduced from u by passing to subspaces is compact and positive.

One has \(\mathrm{Sp}(\widetilde{u})\subset \mathrm{Sp}(u)\) . We also have \(\mathrm{Sp}(\widetilde{u})\subset [0,\lambda_n(u)]\) . Indeed, it suffices to verify that \(\lambda \leqslant \lambda_{n}(u)\) for all \(\lambda \in \mathrm{Sp}_{\mathrm{s}}(\widetilde{u})\) . The number \(\lambda\) is then an eigenvalue of \(u\) , so there exists \(j\in I\) such that \(\lambda = \lambda_j(u)\) . The eigenspace of \(u\) relative to \(\lambda_{j}(u)\) is therefore not contained in F, which implies that \(\lambda_{j}(u)\leqslant \lambda_{n}(u)\) .

Suppose that \(\lambda_n(u) > 0\) . The eigenspace of \(u\) relative to \(\lambda_n(u)\) is then not contained in F, and \(\lambda_n(u)\) therefore belongs to the sensitive spectrum of \(\widetilde{u}\) . We deduce that \(\sup (\mathrm{Sp}(\widetilde{u})) = \lambda_n(u)\) . If \(\lambda_n(u) = 0\) , we have the same result since the spectrum of \(\widetilde{u}\) is then reduced to \(\{0\}\) .

Moreover, we have by definition \(\mathrm{R}_{\mathrm{F}}(u)=\mathrm{R}_{\{0\}}(\widetilde{u})\) , and finally \(\mathrm{R}_{\{0\}}(\widetilde{u})=\sup(\mathrm{Sp}(\widetilde{u}))\) according to prop. 9 of I, p. 139, a). Assertion b) is proved.

PROPOSITION 6. --- For all \(n \in I\) , we have

\[
\lambda_ {n} (u) = \sup _ {\mathrm{F} \in \mathscr {F} _ {n + 1}} r _ {\mathrm{F}} (u) = \inf _ {\mathrm{F} \in \mathscr {F} _ {n}} \mathrm{R} _ {\mathrm{F}} (u).
\]

Let \((e_{n})_{n\in\mathrm{I}}\) be an orthonormal family having the property of prop. 3 of IV, p. 152. For every integer \(n\) such that \(1\leqslant n<M+1\) , let \(F_{n}\in\mathcal{F}_{n}\) be the n-dimensional subspace of E generated by \((e_{0},\ldots,e_{n-1})\) ; by construction, the space \(F_{n}\) is adapted to u. According to prop. 5, we therefore have

\[
\lambda_ {n} (u) = r _ {\mathrm{F} _ {n + 1}} (u) = \mathrm{R} _ {\mathrm{F} _ {n}} (u).\tag{2}
\]

Let \(n \in I\) . Let \(F \in \mathcal{F}_{n+1}\) . The restriction to \(F\) of the orthogonal projector onto \(F_n\) is not injective, so there exists \(x \neq 0\) in \(F\) orthogonal to \(F_n\) .

Since \(x \in \mathrm{F}_{n}^{\circ}\) , we then have (prop. 3 of IV, p. 152)

\[
\begin{array}{l}\langle x\mid u(x)\rangle = \sum_{\substack{m\in \mathrm{I}\\ m\geqslant n}}\lambda_{m}(u)|\langle e_{m}\mid x\rangle |^{2}\\ \leqslant \lambda_{n}(u)\sum_{\substack{m\in \mathrm{I}\\ m\geqslant n}}|\langle e_{m}\mid x\rangle |^{2} = \lambda_{n}(u)\| x\|^{2}. \end{array}
\]

This proves that \(r_{\mathrm{F}}(u)\leqslant\lambda_{n}(u)\) , whence in particular the inequality

\[
\sup _ {\mathrm{F} \in \mathscr {F} _ {n + 1}} r _ {\mathrm{F}} (u) \leqslant \lambda_ {n} (u).\tag{3}
\]

Let \(F \in F_{n}\) . The restriction to \(F_{n+1}\) of the orthogonal projector onto F is not injective, so there exists a vector \(x \neq 0\) in \(F_{n+1}\) orthogonal to F. Since \(x \in F_{n+1}\) , we have (loc. cit.)

\[
\begin{array}{c} \langle x \mid u (x) \rangle = \sum_ {0 \leqslant m \leqslant n} \lambda_ {m} (u) | \langle e _ {m} \mid x \rangle | ^ {2} \\ \geqslant \lambda_ {n} (u) \sum_ {0 \leqslant m \leqslant n} | \langle e _ {m} \mid x \rangle | ^ {2} = \lambda_ {n} (u) \| x \| ^ {2} \end{array}
\]

and therefore \(\mathrm{R_F}(u) \geqslant \lambda_n(u)\) . In particular, it follows that

\[
\inf _ {\mathrm{F} \in \mathcal {F} _ {n}} \mathrm{R} _ {\mathrm{F}} (u) \geqslant \lambda_ {n} (u).\tag{4}
\]

In view of formulas (2), (3) and (4), the proposition is proved.

